\chapter{Implementación y evaluación del prototipo}
\label{cap:implementacion}
Este capítulo documenta cómo se construyó, se desplegó y se verificó la plataforma diseñada en
el Capítulo~\ref{cap:analisis_diseno}. Se organiza siguiendo las capas del sistema: primero
el entorno y las herramientas, después la implementación del lado del servidor, en seguida las
aplicaciones web, luego el empaquetado en contenedores y la publicación en los servidores, y
finalmente la implementación y evaluación de los componentes de inteligencia artificial y las pruebas realizadas. Cierra con un balance
explícito de lo que quedó implementado y lo que no, porque un documento técnico que solo enumera
logros no permite a nadie retomar el trabajo.

\section{Entorno de desarrollo y herramientas}
\label{sec:entorno_herramientas}
La plataforma se construyó con el conjunto de tecnologías cuyos fundamentos se expusieron en la
Sección~\ref{sec:tecnologias}. La Tabla~\ref{tab:herramientas} resume la asignación concreta de
cada herramienta a su función dentro del proyecto.

\rowcolors{2}{rowblue}{white}

\small
\renewcommand{\arraystretch}{1.25}
\setlength{\tabcolsep}{4pt}

\begin{longtable}{|p{3.4cm}|p{3.2cm}|p{7.5cm}|}
	\caption{Herramientas empleadas y su función en el proyecto}
	\label{tab:herramientas} \\
	\hline
	\rowcolor{headerblue}
	\headercell{Capa} & \headercell{Herramienta} & \headercell{Función} \\
	\hline
	\endfirsthead

	\hline
	\rowcolor{headerblue}
	\headercell{Capa} & \headercell{Herramienta} & \headercell{Función} \\
	\hline
	\endhead

	Lenguaje del servidor & Java 21 & Lenguaje de los microservicios de gestión institucional \\ \hline
	Marco del servidor & Spring Boot & Autoconfiguración, servidor embebido, seguridad y
	persistencia \\ \hline
	Construcción & Maven & Gestión de dependencias y empaquetado en archivo ejecutable \\ \hline
	Persistencia & Spring Data JPA & Mapeo objeto-relacional y derivación de consultas \\ \hline
	Base de datos & PostgreSQL 16 & Almacén relacional único compartido \\ \hline
	Autenticación & JSON Web Token & Tokens firmados con algoritmo simétrico \\ \hline
	Documentación de la API & OpenAPI 3 & Especificación e interfaz de exploración por servicio
	\\ \hline
	Hojas de cálculo & Apache POI & Generación del historial exportable y lectura de la
	importación masiva del catálogo \\ \hline
	Cliente web & Angular 21 & Aplicaciones web diferenciadas por rol \\ \hline
	Contenedores & Podman & Empaquetado y ejecución de los servicios y aplicaciones en contenedores \\ \hline
	Servidor web & Nginx & Proxy inverso público y entrega de archivos estáticos \\ \hline
	Certificados & Certbot & Emisión y renovación automática del certificado de transporte
	\\ \hline
	Control de versiones & Git & Historial del proyecto \\ \hline
	Recuperación de antecedentes & Python, scikit-learn y MiniLM & Índices TF-IDF, embeddings semánticos y fusión de rankings \\ \hline
	Resúmenes narrativos & Python, FastAPI y Qwen3 & Selección extractiva y edición institucional \\ \hline
	Inferencia y persistencia & llama.cpp y SQLite & Inferencia cuantizada en CPU y caché persistente de embeddings \\ \hline
\end{longtable}

\normalsize
\setlength{\tabcolsep}{6pt}
\renewcommand{\arraystretch}{1}

Una característica del proyecto que conviene señalar desde el principio: cada microservicio es un
proyecto de construcción independiente, sin un descriptor padre que los agrupe. La decisión
simplifica la evolución de cada servicio por separado, pero tiene una consecuencia operativa
concreta: no existe un comando único que compile todo el sistema, y cada servicio debe compilarse
desde su propio directorio.

\section{Implementación del lado del servidor}
\label{sec:impl_backend}
Los microservicios Java de gestión institucional siguen una estructura interna común. Esta uniformidad
fue una decisión de implementación, no una consecuencia: cuando se construyó cada servicio nuevo
se replicó la organización de carpetas, los filtros y el manejo de errores del anterior, en lugar
de introducir un estilo distinto. El beneficio es que quien conoce un servicio reconoce la estructura de los restantes, y
un defecto corregido en uno se localiza sin dificultad en los demás.

\subsection{Estructura común de un microservicio}
Cada servicio organiza su código en cinco paquetes que corresponden a las capas descritas en la
Sección~\ref{sec:arquitectura_software}:

\begin{itemize}
	\item \texttt{controller}. Puntos de acceso HTTP. Reciben y devuelven objetos de
	transferencia, nunca entidades de persistencia, y declaran el rol requerido para invocarlos.

	\item \texttt{service}. Lógica de negocio y validación de las reglas. Es la única capa que
	conoce las reglas del dominio; los controladores no validan reglas y los repositorios no las
	conocen.

	\item \texttt{repository}. Acceso a datos mediante interfaces de persistencia que derivan las
	consultas a partir del nombre de sus métodos, sin escribir sentencias SQL.

	\item \texttt{entity}. Entidades de persistencia con su correspondencia a tablas y columnas.

	\item \texttt{model}. Objetos de transferencia de datos para las peticiones y las respuestas.
\end{itemize}

La separación entre \texttt{entity} y \texttt{model} merece una justificación explícita, porque
implica escribir clases que a primera vista parecen redundantes. Exponer directamente la entidad
de persistencia en una respuesta HTTP acopla el contrato público del servicio a la estructura de
la base de datos: cualquier columna nueva aparece automáticamente en la respuesta, incluidas las
que no deberían ser públicas ---la credencial cifrada de un usuario, por ejemplo---. Mantener
objetos de transferencia separados obliga a decidir de forma consciente qué se expone.

\subsection{Verificación de tokens y control de acceso}
El mecanismo de autenticación se implementó como un filtro que se ejecuta antes de que la
petición llegue al controlador. Su lógica es compartida por los servicios Java que requieren autenticación: extrae el token del
encabezado de autorización, verifica su firma con la clave compartida, extrae de él la identidad
y el rol, y establece el contexto de seguridad de la petición. Si el token falta, está vencido o
su firma no coincide, el filtro rechaza la petición sin llegar a la lógica de negocio.

Sobre ese contexto se aplica la autorización, declarada en cada controlador mediante anotaciones
que exigen un rol determinado. Un detalle de configuración que resultó relevante durante la
implementación: cuando un mismo prefijo de ruta tiene una parte pública y una parte
administrativa, el orden de declaración de las reglas importa, porque el marco de seguridad
aplica la primera coincidencia. Declarar la regla pública antes de la administrativa hace que la
menos restrictiva prevalezca por ser más genérica, y el resultado es un punto de acceso
administrativo efectivamente abierto. Esta situación se atendió declarando siempre la regla más
específica primero, y quedó registrada como regla de negocio RN-A05 justamente para que no se
reintroduzca al agregar servicios nuevos.

\subsection{Manejo centralizado de errores}
Cada servicio implementa un manejador global de excepciones que traduce cualquier error de la
capa de negocio a una respuesta con estructura uniforme: mensaje, marca de tiempo y código. La
alternativa ---capturar excepciones en cada controlador--- se descartó porque produce respuestas
de error con formatos distintos según el punto de acceso, y obliga al cliente a implementar
tantas variantes de manejo de error como puntos de acceso consuma. Con el manejador centralizado,
las aplicaciones web leen el mensaje de error siempre del mismo lugar.

\subsection{Documentación de la interfaz de programación}
Los servicios Java exponen su especificación en formato OpenAPI y una interfaz de exploración
interactiva. Esta decisión resolvió un problema concreto de coordinación: durante el desarrollo
del cliente, la única forma de conocer la forma exacta de una petición era leer el código del
controlador correspondiente. Con la especificación publicada, el contrato de cada servicio es
consultable sin abrir el código.

\subsection{Puntos de acceso implementados}
La Tabla~\ref{tab:endpoints} presenta los puntos de acceso principales agrupados por servicio,
con su condición de acceso. La especificación completa de cada uno, con sus parámetros y
respuestas, se consulta en la interfaz de exploración de cada servicio y se reproduce en el anexo
correspondiente.

\rowcolors{2}{rowblue}{white}

\footnotesize
\renewcommand{\arraystretch}{1.2}
\setlength{\tabcolsep}{4pt}

\begin{longtable}{|p{2.6cm}|p{8.6cm}|p{3.0cm}|}
	\caption{Puntos de acceso principales por microservicio}
	\label{tab:endpoints} \\
	\hline
	\rowcolor{headerblue}
	\headercell{Servicio} & \headercell{Punto de acceso} & \headercell{Acceso} \\
	\hline
	\endfirsthead

	\hline
	\rowcolor{headerblue}
	\headercell{Servicio} & \headercell{Punto de acceso} & \headercell{Acceso} \\
	\hline
	\endhead

	\texttt{auth-service} & \texttt{POST /api/auth/login} & Público \\ \hline
	& \texttt{POST /api/auth/activar-cuenta} & Público \\ \hline
	& \texttt{POST /api/auth/solicitar-codigo} & Público \\ \hline
	& \texttt{POST /api/auth/reset-password} & Público \\ \hline
	& \texttt{GET} y \texttt{PUT /api/auth/me}, \texttt{/api/auth/perfil} & Sesión de quejoso
	\\ \hline

	\texttt{queja-service} & \texttt{POST /api/quejoso/\allowbreak quejas/registro-publico} & Público
	\\ \hline
	& \texttt{GET /api/quejoso/quejas/folio/\{folio\}} & Público \\ \hline
	& \texttt{POST /api/quejoso/quejas/registrar} & Sesión de quejoso \\ \hline
	& \texttt{GET /api/quejoso/quejas/mias} & Sesión de quejoso \\ \hline
	& \texttt{GET} y \texttt{PUT /api/quejoso/\allowbreak quejas/mias/\{folio\}} & Sesión de quejoso \\ \hline
	& \texttt{GET} y \texttt{PUT /api/quejoso/conciliaciones/**} & Sesión de quejoso \\ \hline

	\texttt{notificaciones-\allowbreak service} & \texttt{GET /api/notificaciones/mias} y
	\texttt{/mias/no-leidas} & Sesión de quejoso \\ \hline
	& \texttt{PUT /api/notificaciones/\{id\}/leida} & Sesión de quejoso \\ \hline
	& \texttt{POST /api/notificaciones/enviar} y \texttt{/registrar} & Uso interno entre
	servicios \\ \hline

	\texttt{catalogo-\allowbreak service} & \texttt{GET /api/catalogos/dependencias} y
	\texttt{/dependencias/\{clave\}} & Público \\ \hline
	& \texttt{GET}, \texttt{POST}, \texttt{PUT /api/catalogos/\allowbreak dependencias/admin/**} &
	Administración \\ \hline
	& \texttt{POST /api/catalogos/\allowbreak dependencias/admin/\allowbreak importar-excel} & Administración \\ \hline

	\texttt{admin-service} & \texttt{POST /api/admin/auth/login} & Público \\ \hline
	& \texttt{GET}, \texttt{POST}, \texttt{PUT}, \texttt{DELETE /api/admin/personal/**} &
	Administración \\ \hline
	& \texttt{GET} y \texttt{PUT /api/admin/plantillas/**} & Administración \\ \hline
	& \texttt{GET} y \texttt{POST /api/admin/\allowbreak seguridad/respaldos/**} & Administración \\ \hline
	& \texttt{POST /api/admin/seguridad/restaurar} & Administración \\ \hline
	& \texttt{GET /api/admin/seguridad/bitacora} & Administración \\ \hline
	& \texttt{GET /api/admin/dashboard/resumen} & Administración \\ \hline
	& \texttt{PUT /api/admin/perfil/password} & Cualquier cuenta de personal \\ \hline

	\texttt{revision-\allowbreak service} & \texttt{GET /api/revision/bandeja} & Recepción \\ \hline
	& \texttt{GET /api/revision/quejas/\{folio\}} y \texttt{/antecedentes} & Recepción \\ \hline
	& \texttt{GET /api/revision/\allowbreak quejas/evidencias/\{id\}} & Recepción \\ \hline
	& \texttt{POST /api/revision/\allowbreak quejas/\{folio\}/rechazar} y \texttt{/turnar} & Recepción
	\\ \hline
	& \texttt{POST /api/revision/registro-manual} & Recepción \\ \hline
	& \texttt{GET /api/revision/historial} y \texttt{/historial/exportar} & Recepción \\ \hline
	& \texttt{GET /api/revision/catalogos/areas} y \texttt{/defensores} & Recepción \\ \hline
	& \texttt{POST} y \texttt{GET /api/revision/conciliaciones} & Recepción \\ \hline

	\texttt{chatbot-\allowbreak service} & \texttt{GET /api/chatbot/menu} & Público \\ \hline
	& \texttt{GET}, \texttt{POST}, \texttt{PUT}, \texttt{DELETE /api/chatbot/\allowbreak admin/preguntas/**} &
	Administración \\ \hline
	\texttt{antecedentes-service} & \texttt{POST /api/antecedentes/buscar} & Uso interno \\ \hline
	\texttt{resumen-service} & \texttt{POST /resumir} & Uso interno \\ \hline
\end{longtable}

\normalsize
\setlength{\tabcolsep}{6pt}
\renewcommand{\arraystretch}{1}

\section{Implementación de las aplicaciones web}
\label{sec:impl_frontend}
Las aplicaciones web se desarrollaron como aplicaciones de página única con componentes
autónomos, sin el sistema de módulos de versiones anteriores del marco de trabajo. Cada
aplicación es un proyecto independiente con su propio árbol de rutas, su guarda de
autenticación y sus interceptores.

\subsection{Estructura común}
Estas aplicaciones comparten una organización general: una capa \texttt{core} con los servicios de
comunicación con el servidor, los modelos de datos y la configuración de seguridad del cliente;
una capa \texttt{shared} con los componentes reutilizables ---encabezado, pie, avisos, ayuda
flotante---; y una capa \texttt{pages} con un componente por pantalla, cargado de forma diferida
para que el navegador no descargue el código de pantallas que el usuario no ha visitado.

\subsection{Interceptores y guardas}
Dos piezas transversales resuelven problemas que de otro modo habría que repetir en cada
pantalla.

El \textbf{interceptor de autorización} agrega el token de sesión al encabezado de cada petición
saliente, de modo que ningún servicio de la aplicación tiene que recordar hacerlo. El mismo
interceptor detecta las respuestas de sesión inválida o insuficiente y, en ese caso, cierra la
sesión local y redirige al inicio de sesión. Esto implementa el requerimiento RNF-Q06 en un solo
punto del código.

La \textbf{guarda de rutas} impide el acceso a las secciones autenticadas cuando no existe una
sesión válida en el cliente. Es importante entender su alcance: la guarda es una medida de
usabilidad, no de seguridad, porque se ejecuta en el navegador y puede eludirse. La seguridad
real la aplica el servidor, que rechaza cualquier petición sin token válido con independencia de
lo que haga el cliente. Confundir ambas cosas es un error frecuente y por eso se hace explícito.

\subsection{Detección de cambios sin zonas}
La versión del marco de trabajo empleada permite ejecutar la aplicación sin la biblioteca que
tradicionalmente intercepta las operaciones asíncronas para disparar la actualización de la
vista. Ese modo reduce el trabajo del navegador, pero tiene una consecuencia que apareció durante
la implementación: la vista no se actualizaba de forma automática al recibir la respuesta de una
petición, de modo que los datos llegaban pero la pantalla seguía mostrando el estado anterior.

La solución adoptada fue un interceptor adicional que solicita de forma explícita una
actualización de la vista tras cada respuesta del servidor. Se documenta aquí porque es una
particularidad no evidente del entorno y porque cualquier pantalla nueva que se agregue al
proyecto depende de que ese interceptor siga registrado.

\section{Empaquetado en contenedores y despliegue}
\label{sec:impl_despliegue}
El sistema se ejecuta mediante contenedores distribuidos según sus responsabilidades, sobre la infraestructura descrita en la
Sección~\ref{sec:diseno_arquitectura}. La Tabla~\ref{tab:contenedores} presenta los principales componentes de despliegue, sin incluir los contenedores auxiliares de prueba.

\rowcolors{2}{rowblue}{white}

\small
\renewcommand{\arraystretch}{1.2}
\setlength{\tabcolsep}{4pt}

\begin{longtable}{|p{3.4cm}|p{3.0cm}|p{7.6cm}|}
	\caption{Matriz de despliegue en contenedores}
	\label{tab:contenedores} \\
	\hline
	\rowcolor{headerblue}
	\headercell{Contenedor} & \headercell{Servidor} & \headercell{Contenido} \\
	\hline
	\endfirsthead

	\hline
	\rowcolor{headerblue}
	\headercell{Contenedor} & \headercell{Servidor} & \headercell{Contenido} \\
	\hline
	\endhead

	Base de datos & Aplicación & PostgreSQL 16 con almacenamiento persistente \\ \hline
	\texttt{auth-service} & Aplicación & Imagen dedicada del servicio de autenticación \\ \hline
	\texttt{quejas-service} & Aplicación & Imagen dedicada del servicio de quejas \\ \hline
	\texttt{notificaciones-\allowbreak service} & Aplicación & Imagen dedicada del servicio de
	notificaciones \\ \hline
	\texttt{catalogo-\allowbreak service} & Aplicación & Imagen dedicada del servicio de catálogo \\ \hline
	\texttt{admin-service} & Aplicación & Imagen propia que incluye el cliente de la base de
	datos, necesario para generar y restaurar respaldos \\ \hline
	\texttt{revision-\allowbreak service} & Aplicación & Imagen dedicada del servicio de revisión \\ \hline
	\texttt{chatbot-\allowbreak service} & Aplicación & Imagen dedicada del servicio de preguntas frecuentes
	\\ \hline
	Portal del quejoso & Acceso público & Servidor de estáticos con la aplicación compilada
	\\ \hline
	Consola de administración & Acceso público & Servidor de estáticos con la aplicación compilada
	\\ \hline
	Panel de revisión & Acceso público & Servidor de estáticos con la aplicación compilada
	\\ \hline
	Proxy inverso & Acceso público & Nginx con la configuración de enrutamiento y el certificado
	de transporte \\ \hline
	Servicio de antecedentes & Aplicación / IA & Recuperación léxica y semántica y caché de embeddings \\ \hline
	Servicio de resúmenes & Aplicación / IA & Selección extractiva y edición institucional \\ \hline
	Servidor Qwen3 & Aplicación / IA & Inferencia cuantizada en CPU con llama.cpp \\ \hline
	Panel de Primer Contacto & Acceso público & Interfaz de análisis asistido por IA \\ \hline
	Panel de Subdefensoría & Acceso público & Interfaz institucional de gestión de expedientes \\ \hline
\end{longtable}

\normalsize
\setlength{\tabcolsep}{6pt}
\renewcommand{\arraystretch}{1}

\subsection{Configuración externalizada}
Cada microservicio monta su configuración de producción como un volumen del contenedor, en lugar
de incluirla en la imagen. La consecuencia práctica es una distinción importante en el
procedimiento de operación: un cambio de configuración ---una dirección, un tiempo de
expiración--- se aplica editando el archivo montado y reiniciando el contenedor, mientras que un
cambio de lógica exige recompilar el servicio, generar un artefacto nuevo y reconstruir la
imagen. Esta separación es la que permite que la misma imagen sirva en cualquier ambiente sin
reconstruirse, tal como se expuso en la Sección~\ref{sec:tecnologias}.

Durante la operación apareció una consecuencia no evidente de este esquema: editar el archivo de
configuración en el servidor anfitrión no basta para que el contenedor lo tome, porque el punto de
montaje conserva la referencia al archivo original. Es necesario reiniciar el contenedor, no solo
recargar el servicio. Se documenta porque es una fuente de diagnósticos erróneos: el archivo se ve
correcto en el servidor y el servicio sigue comportándose como antes.

\subsection{Publicación segura y cifrado del canal}
La publicación al exterior se resolvió en tres pasos. Primero se restringió el cortafuegos del
servidor de aplicación para que los puertos de los microservicios y de la base de datos aceptaran
conexiones únicamente desde la dirección del servidor de acceso público, dejando fuera cualquier
otro origen de Internet. Después se configuró el proxy inverso del servidor de acceso público con
las reglas de enrutamiento de la Tabla~\ref{tab:enrutamiento}. Por último se emitió el certificado
de transporte mediante el método de validación por directorio web, se reconfiguró el proxy para
atender el puerto cifrado y se dejó el puerto sin cifrar únicamente para redirigir el tráfico y
atender la renovación automática del certificado.

Durante este proceso se documentaron dos incidentes que conviene registrar por su valor
diagnóstico. El primero fue una aparente imposibilidad de alcanzar los puertos web del servidor
de acceso público desde el exterior, mientras que el puerto de administración remota respondía
con normalidad. Tras descartar de forma sistemática la configuración del servicio, la publicación
del contenedor y el cortafuegos local, se identificó que la causa era una regla de cortafuegos
del proveedor que estaba guardada pero no sincronizada con la máquina virtual: existía en el panel
de control pero no se había aplicado. El segundo fue un aviso recurrente del motor de contenedores
sobre una red virtual mal configurada, que provenía de un intento previo de crear una red propia
que finalmente no se utilizó; su eliminación resolvió el aviso. Ambos casos ilustran un principio
de diagnóstico útil: cuando un componente responde de forma local y no desde el exterior, el
problema no está en el componente sino en el camino.

\section{Implementación de la búsqueda de antecedentes}
\label{sec:impl_antecedentes}
El servicio de búsqueda de antecedentes se implementó en Python y se expone mediante una
interfaz HTTP. Para cada consulta recibe la descripción de una queja, recupera registros
existentes y devuelve documentos ordenados para su revisión por el personal autorizado.
La recuperación combina dos representaciones complementarias. La primera utiliza TF-IDF
con unigramas y bigramas para identificar coincidencias léxicas; la segunda obtiene
representaciones semánticas con un modelo multilingüe de la familia MiniLM, a través
de Sentence Transformers. Las puntuaciones de ambas técnicas no se suman directamente:
cada recuperador genera su propia clasificación de candidatos y estos ordenamientos
se integran mediante Reciprocal Rank Fusion (RRF) ponderado, descrito en el Marco Teórico.
En la configuración final adoptada para el sistema se asignó un peso del 75\,\% al
ordenamiento de TF-IDF y del 25\,\% al de MiniLM. Estos pesos corresponden
a la fusión de rankings, no a una suma directa de las similitudes. Su elección
representa una decisión de implementación y no demuestra una configuración
universalmente óptima para cualquier conjunto de quejas.

\subsection{Preparación del índice y fragmentación}
Los documentos se leen desde las tablas de quejas e histórico mediante consultas de
solo lectura. Para la comparación semántica, el contenido extenso se divide en
fragmentos de hasta 128 tokens, con un solapamiento de 24 tokens. Esta fragmentación
permite representar secciones diferentes de un mismo expediente, en lugar de depender
únicamente de una representación global de toda la narración.

El procesamiento inicial obtiene las representaciones vectoriales y las almacena en
una caché persistente basada en SQLite. La clave de cada entrada depende de una
huella del contenido y de la configuración de representación. Cuando se actualiza
el índice, las entradas de documentos sin cambios se reutilizan; las nuevas o
modificadas se procesan de nuevo. La construcción ocurre sin bloquear el inicio
de la API y el índice vigente se sustituye cuando la actualización termina.

\subsection{Recuperación y ordenamiento}
La consulta se transforma tanto en un vector TF-IDF como en representaciones
semánticas. Para interpretar y reproducir el resultado de documentos extensos,
es necesario registrar cómo se agregan las puntuaciones de sus fragmentos y cómo
se excluyen, cuando corresponde, las autocoincidencias de las consultas de evaluación.
La información experimental presentada no permite verificar por sí sola estos dos aspectos. Se generan listas de candidatos de cada recuperador y se combinan sus
posiciones mediante RRF ponderado. La implementación emplea una constante de
suavizado de 60 y ponderaciones finales 0.75 para TF-IDF y 0.25 para MiniLM.
Las demás combinaciones se examinaron experimentalmente; no se presupone que
la configuración desplegada sea óptima para cualquier tipo de consulta.

El servicio conserva las medidas léxica y semántica por separado, además de la
puntuación de fusión. La puntuación RRF es una señal de ordenamiento, no una
probabilidad de que dos quejas se refieran al mismo hecho. Por ello, cualquier porcentaje
mostrado en la interfaz debe identificarse por su origen (por ejemplo, similitud léxica
o semántica), sin presentarlo como una probabilidad calibrada de duplicidad. Los resultados pueden
agruparse por semejanzas narrativas o coincidencias de personas y otros datos,
pero la confirmación de duplicidad permanece bajo revisión humana.

\section{Implementación de la generación de resúmenes}
\label{sec:impl_resumenes}
El servicio de resúmenes utiliza una secuencia extractiva-generativa. Primero,
un procedimiento de selección documental basado en relevancia, redundancia y
reglas de conservación de contexto obtiene una captura extractiva a partir de
la narración de la queja. La selección incorpora el principio de Maximum
Marginal Relevance (MMR) y criterios adicionales para priorizar hechos y
reducir repeticiones. Posteriormente, un modelo generativo basado en Qwen3 recibe la captura
seleccionada y produce un borrador de estilo institucional y en tercera
persona, utilizando instrucciones para conservar actores, negaciones,
fechas, cantidades, evidencias y solicitudes. La denominación de editor
lingüístico describe su función prevista, pero no garantiza que el modelo
conserve todos los hechos de la captura.

\subsection{Ejecución del editor lingüístico}
El editor utiliza un modelo Qwen3 de cuatro mil millones de parámetros,
cuantizado en formato Q4\_K\_M y servido por llama.cpp. La inferencia se ejecuta
en CPU mediante un proceso independiente del servicio HTTP que expone el
resumen. Esta separación permite modificar el procesamiento extractivo sin
reconstruir el servidor del modelo y viceversa.

La captura se divide en bloques de hasta seis oraciones y con una referencia
de 120 palabras como límite por bloque. Cada bloque se edita secuencialmente,
con un límite de generación de hasta 450 tokens. Se implementó un presupuesto
aproximado de 95 segundos para la etapa de edición, con el objetivo de reducir
el riesgo de agotar el tiempo de espera del consumidor Java. Cuando un bloque
no puede editarse satisfactoriamente, se conserva su contenido extractivo
original. Este mecanismo evita perder el bloque por un fallo de generación,
pero no garantiza la fidelidad completa de un resumen.

\subsection{Respuesta y revisión humana}
El servicio expone una operación HTTP para generar el borrador y devuelve tanto
el texto editado como la captura extractiva de la que partió. La respuesta
indica que requiere revisión humana. El resultado no constituye una resolución
jurídica ni reemplaza la queja original; su propósito es apoyar la lectura
inicial del expediente.


\section{Pruebas y verificación}
\label{sec:pruebas}
La verificación de la plataforma combinó pruebas automatizadas de lógica de negocio,
comprobaciones de integración, revisiones de configuración y recorridos funcionales
manuales. Estas actividades ofrecen evidencias de distinto alcance y no deben
interpretarse como equivalentes. Los casos, sus resultados y las verificaciones
pendientes se detallan en el Capítulo~\ref{cap:pruebas}.

\begin{itemize}
    \item \textbf{Compilación y empaquetado.} Los microservicios Java se construyeron
    con Maven y las aplicaciones web se compilaron en modo de producción. Esta
    verificación identifica errores de construcción, pero no acredita por sí sola
    el cumplimiento de las reglas de negocio.
    \item \textbf{Pruebas unitarias automatizadas.} Se utilizaron JUnit~5 y Mockito
    para verificar lógica y validaciones de los microservicios sin depender de la
    base de datos ni de servicios externos. En la ejecución documentada se reportaron
    364 casos aprobados, sin fallas ni errores. Posteriormente se incorporaron
    ocho casos de resiliencia, cuyo resultado debe distinguirse del conteo inicial.
    \item \textbf{Integración y puntos de acceso.} Se realizaron comprobaciones
    de contratos HTTP, autorización, interacción entre servicios y conectividad a
    través del proxy inverso. Algunas fueron ejecutadas dinámicamente, otras se
    comprobaron mediante inspección de código o configuración y otras permanecen
    pendientes; el Capítulo~\ref{cap:pruebas} especifica su estado por caso.
    \item \textbf{Recorridos de sistema.} Se revisaron manualmente flujos de
    registro, recepción, turnado y seguimiento. La automatización integral de los
    recorridos de extremo a extremo mediante Playwright aún no estaba concluida.
\end{itemize}

En consecuencia, no es correcto afirmar que toda la verificación fue manual ni
que existe una cobertura automatizada integral. Las pruebas unitarias proporcionan
una base de regresión para reglas específicas; la cobertura de líneas, los casos
pendientes de seguridad, capacidad, aceptación y los flujos completos todavía
requieren evidencias adicionales. Tampoco se documenta aquí la existencia de una
integración continua que ejecute automáticamente todas las suites.

